\label{sec:limitations}
\paragraph{Evidence.} Our results rest on 17 benchmarks from the competition's public data, with random label acquisition, and three benchmarks serve as test. Intervals resample tasks, so they do not reflect which benchmarks are held out. We therefore also report the gain over all 680 choices of three held-out benchmarks (Appendix~\ref{app:split-sensitivity}), and read them as a measure of sensitivity, since they overlap.

\paragraph{Selection.} All settings were chosen on validation. Test scores were visible during development, so we treat the test split as one of many possible splits and report the others. Our account of when the prior helps was formed after seeing the results and uses observed outcomes, so it explains them and does not select a prior for a new benchmark.

\paragraph{Not tested.} Sharing labels across subjects, other label-selection policies, priors from subject metadata, reasoning-trace features (unavailable at test time) and stronger difficulty predictors. The conclusions are limited to the models and data above.

